\section{Introduction}
\label{sec:introduction}

Rotating machinery is central to manufacturing, transportation, energy, and other asset-intensive industrial systems, where unexpected failures can interrupt production, increase maintenance expenditure, and compromise operational safety. Predictive and condition-based maintenance seek to reduce these risks by using condition-monitoring information to support timely diagnostic and maintenance decisions \citep{Jardine2006CBM,Nunes2023Challenges}. Among the available condition-monitoring modalities, vibration signals remain particularly important because bearing, rotor, and other mechanical defects modify the temporal and spectral structure of the measured response. The increasing availability of industrial sensing and computing has consequently driven a transition from manually engineered diagnostic features toward data-driven and deep-learning methods for machine health monitoring \citep{Zhao2019DLHealth,Souza2021RotatingDL,Nascimento2022T4PdM}.

Deep neural networks can learn discriminative fault representations directly from measured signals, but many diagnostic models remain closely tied to a particular dataset, machine configuration, or label space. This limits reuse when fault classes, sampling conditions, and acquisition systems change. Recent research has therefore moved beyond purely task-specific classifiers toward reusable pretrained representations. FaultFormer demonstrated self-supervised masked pretraining for bearing vibration signals and subsequent adaptation to low-data and cross-dataset settings \citep{Zhou2024FaultFormer}; BearingFM proposed a foundation-style bearing model using domain knowledge and contrastive learning \citep{Lai2024BearingFM}; and RmGPT introduced a unified pretrained framework for rotating-machinery diagnosis and prognosis \citep{Wang2025RmGPT}. More recently, VibFM scaled masked spectrogram modelling to 16 open vibration datasets and evaluated transfer to held-out bearing datasets \citep{Mannone2026VibFM}. These developments indicate a clear shift toward representations that can support multiple diagnostic tasks rather than a separate network for every individual dataset.

At the same time, sequence modelling for machinery diagnosis has begun to adopt state-space architectures. Mamba introduced selective state-space modelling as an alternative to attention-based sequence processing, with input-dependent state updates and linear sequence scaling \citep{Gu2023Mamba}. In rotating-machinery applications, MD-BiMamba employed bidirectional Mamba processing for inter-shaft bearing diagnosis \citep{Wang2025MDBiMamba}, while VibrMamba combined bidirectional state-space processing with lightweight projection layers and evaluated six classification tasks across four public machinery datasets \citep{Yi2025VibrMamba}. These studies show that Mamba-based and lightweight Mamba-based fault diagnosis are already established research directions. Consequently, the relevant deployment question is no longer whether a compact Mamba classifier can be designed, but whether an already trained reusable state-space representation can be compressed in a controlled manner without introducing an unacceptable loss of diagnostic capability.

This question is important because reusable encoders can still impose substantial computational and storage costs at inference time. In industrial deployment, model quality must be balanced against memory footprint, arithmetic complexity, latency, throughput, and the capabilities of the target hardware. Knowledge distillation and structural compression have therefore received increasing attention in bearing diagnosis. DTCP-MAKD combines decreasing-threshold channel pruning with multi-assistant knowledge distillation to bridge the capacity gap between teacher and student models \citep{Zhong2024DTCPMAKD}; BFD-KD uses multiple teachers with dynamic weighting and intermediate-feature fusion for lightweight rolling-bearing diagnosis \citep{Sun2026BFDKD}; and BearingPGA-Net combines decoupled knowledge distillation, customized quantization, and FPGA acceleration for deployable bearing diagnosis \citep{Liao2024BearingPGANet}. These works demonstrate that structural reduction, multi-teacher supervision, quantization, and hardware-oriented deployment are not individually new. However, they primarily compress task-specific diagnostic networks or directly design lightweight architectures rather than examining the compression of a shared multi-dataset selective-state-space encoder.

The distinction is relevant for the present study. The full Physical-Coordinate Spectro-Temporal Encoder (PC-STE) uses one shared encoder across CWRU, JNU, HIT, and MaFaulDa, with dataset-specific classification heads. Its Full-S1 configuration is initialized through masked spectrogram reconstruction and employs four bidirectional temporal Mamba blocks together with physical time--frequency coordinates and cross-band mixing. The purpose of the present work is not to re-establish self-supervised pretraining, Mamba-based fault diagnosis, or multi-teacher knowledge distillation as new concepts. Instead, it addresses a narrower deployment problem: \emph{can a shared multi-dataset bidirectional state-space encoder be structurally compressed while retaining its diagnostic performance and delivering measurable computational benefits?}

To investigate this question, the bidirectional Full-S1 encoder is structurally reduced to K1, a four-layer forward-only student that removes the reverse temporal scan while preserving the remaining encoder hierarchy and the shared multi-dataset role. K1 is then trained by knowledge distillation from an ensemble of three independently seeded Full-S1 teachers from the same global fold. This teacher construction differs from heterogeneous multi-teacher or pruning-derived teacher-assistant schemes: all teachers share the same architecture, fold, task definition, and training protocol, with their diversity arising primarily from stochastic optimization. The resulting student is evaluated as a compression of the existing encoder rather than as an unrelated small model designed from scratch.

A central methodological feature of the study is that compression is treated as a \emph{performance-retention} problem. The final Full-S1--K1 comparison uses three global folds and three matched random seeds, yielding nine paired experimental cells. Before final TEST evaluation, a Macro-F1 non-inferiority margin of $-0.02$ was fixed to define the largest acceptable loss in the equal-domain Macro-4 endpoint. This design permits the efficiency gain to be assessed against an explicit retention criterion rather than through a single accuracy comparison or an informal statement of ``minimal degradation.'' Computational evaluation is performed separately using fixed validation inputs and a common hardware environment, with parameter count, serialized size, FLOPs, corrected CPU/GPU latency, throughput, and GPU memory reported under controlled conditions. A post-training Q8 stage is additionally examined as a secondary deployment extension and is kept distinct from the primary Full-S1--K1 confirmatory study.

The final results show that structural compression does not require a loss of aggregate diagnostic performance in this setting. Across the nine matched cells, Macro-4 Macro-F1 increases from $0.934644\pm0.023258$ for Full-S1 to $0.955334\pm0.018393$ for K1, corresponding to a paired change of $+0.020691\pm0.013509$, with K1 higher in eight of the nine cells. The predefined non-inferiority criterion is satisfied by the exact paired sign-flip test ($p=0.001953125$). At the same time, K1 reduces encoder parameters by $42.24\%$, serialized FP32 model size by $42.14\%$, and equal-domain computational cost by $45.38\%$. Under the corrected inference benchmark, K1 provides a $1.874\times$ end-to-end CPU batch-1 speed-up, a $1.633\times$ end-to-end GPU batch-1 speed-up, and approximately $2.0\times$ GPU batch-32 throughput. The secondary Q8 representation retains essentially the same aggregate Macro-F1 while reducing stored model size by a further $71.12\%$ relative to K1 FP32; because the preserved implementation has no true INT8 CUDA path, Q8 is interpreted primarily as a storage and CPU-deployment extension rather than a GPU acceleration result. The study therefore focuses on multi-dataset performance retention and deployment efficiency, not on claims of unseen-machine generalization.

\paragraph{Contributions.}
The main contributions of this study are summarized as follows:
\begin{enumerate}[leftmargin=*,label=(\roman*)]
    \item \textbf{Directionality-specific structural compression of a shared PC-STE encoder.}
    The bidirectional Full-S1 selective-state-space encoder is structurally reduced to a forward-only four-layer K1 student while retaining one shared encoder across CWRU, JNU, HIT, and MaFaulDa with dataset-specific classification heads.

    \item \textbf{Ensemble knowledge distillation for the compressed encoder.}
    K1 is distilled from same-fold, independently seeded Full-S1 teachers, combining classification distillation with relational supervision while preserving a single compact inference model.

    \item \textbf{Matched performance-retention evaluation.}
    The teacher--student comparison uses three global folds and three random seeds, producing nine matched cells, together with a prespecified $-0.02$ Macro-F1 non-inferiority margin to assess whether compression introduces an unacceptable diagnostic-performance loss.

    \item \textbf{Controlled predictive--efficiency characterization with a secondary Q8 extension.}
    The study reports model size, parameter count, FLOPs, corrected CPU/GPU latency, throughput, and memory under one controlled benchmark environment, and additionally evaluates post-training Q8 quantization as a secondary deployment extension.
\end{enumerate}

The remainder of this paper is organized as follows. Section~\ref{sec:literature} reviews self-supervised and foundation-style vibration models, Mamba-based machinery diagnosis, lightweight fault-diagnosis approaches, and deployment-oriented compression. Section~\ref{sec:methodology} presents the datasets, PC-STE representation, Full-S1 reference model, K1 structural compression, ensemble distillation, Q8 extension, and evaluation protocol. Section~\ref{sec:results} reports predictive performance, non-inferiority analysis, computational efficiency, and Q8 results. Finally, Section~\ref{sec:conclusion} summarizes the main findings, limitations, and future work.
